\chapter{Grids and Fourier transforms}\label{ch:grids-fft}
\chapterstatus{written}{2026-11-14}

Fields are stored on distributed grids with in-place FFTW transforms \cite{Frigo2005}: real grids with real-to-complex
transforms for densities and potentials, complex grids for wave functions and the NUFFT, and generic MPI grids of any type
for the multigrid solvers. Generic algorithms are written against one interface, so they run on slabs and on pencils.

\begin{modulebox}{FFTWGrid}
\nsp{FML::GRID}: \file{FFTWGrid/FFTWGrid.h} (\cls{FFTWGrid}), \file{FFTWGrid/CFFTWGrid.h} (\cls{CFFTWGrid}),
\file{FFTWGrid/FFTWGlobal.h} (precision), \file{FFTWGrid/GridInterface.h}, \file{FFTWGrid/GhostedCells.h}
(\cls{GhostedCells}), \file{FFTWGrid/FourierResampling.h} (\code{resize_fourier_grid}, \cls{FourierResamplingPlan}),
\file{FFTWGrid/CopyFourierModes.h}.\par
Documentation: \file{FFTWGrid/FourierResampling.md}. Tests: \file{tests/test_grids.cpp}, \file{tests/test_complex_grids.cpp},
\file{tests/test_grid_generic.cpp}, \file{tests/test_ghosted_cells.cpp}, \file{tests/test_fourier_resampling.cpp}.
\end{modulebox}
\begin{modulebox}{MPIGrid}
\nsp{FML::GRID}: \file{MPIGrid/MPIGrid.h} (\cls{MPIGrid}), \file{MPIGrid/MPIMultiGrid.h} (\cls{MPIMultiGrid}),
\file{MPIGrid/ConvertMPIGridFFTWGrid.h}. The adaptive grid \cls{AdaptiveMPIGrid} is described with its solver in
\cref{sec:mg-amr}.
\end{modulebox}

\section{FFTWGrid}\label{sec:grids-fftwgrid}

\cls{FFTWGrid} holds an $N_{\rm mesh}^D$ real field ($D\ge2$) split in x-slabs, with the real and the Fourier grid in the
same array: the transforms are in place. The real grid is padded to $2(N_{\rm mesh}/2+1)$ along the last axis; the Fourier
grid has $N_{\rm mesh}^{D-1}(N_{\rm mesh}/2+1)$ complex cells, local in $k_x$. The memory is
\[
[\text{extra slices left}]\;[\text{main grid}]\;[\text{extra slices right}],
\]
so a cell $k$ slices to the left is at \code{index - k * stride_x}, also into the extra slices. The grid records whether it
holds real or Fourier data.

\paragraph*{Conventions.} \code{fftw_r2c()} divides by $N_{\rm mesh}^D$, so the Fourier values are the amplitudes of the
discrete Fourier series, $\delta(\bm x)=\sum_{\bm k}\delta_{\bm k}e^{i\bm k\cdot\bm x}$, and \code{fftw_c2r()} does not
normalise (\cref{ch:conventions}). Wave-vectors are $\bm k=2\pi\bm n$ in box units (multiply by $1/B$ for physical units),
with $n_i\in(-N/2,N/2]$ and $n_D\in[0,N/2]$. A mode with $0<n_D<N/2$ also stands for $-\bm k$ (weight 2 in sums over the
grid; \code{fourier_mode_multiplicity}). The DC mode lives on the task whose slab starts at $x=0$ (\code{has_dc_mode}).
The precision is double by default, float with \code{SINGLE_PRECISION_FFTW} and long double with
\code{LONG_DOUBLE_PRECISION_FFTW}.

\paragraph*{Access.} By local integer coordinates or by index (\code{get_real}, \code{set_real}, \code{add_real},
\code{get_fourier}, \ldots, with \code{_from_index} variants), positions in $[0,1)^D$ (\code{get_real_position}),
wave-vectors and their norms by index, filling from a function of position or wave-vector, and range iterators over the real
or Fourier cells, optionally restricted to rows for OpenMP:

\begin{cppcode}[Filtering a field in Fourier space]
#include <FML/FFTWGrid/FFTWGrid.h>
FML::GRID::FFTWGrid<3> grid(Nmesh, 1, 1);            // one extra slice on each side
grid.fill_real_grid([](std::array<double, 3> & x) { return f(x); });
grid.fftw_r2c();                                     // normalised by 1 / Nmesh^3
#pragma omp parallel for
for (IndexIntType irow = 0; irow < grid.get_local_nrows(); irow++)
    for (auto && i : grid.get_fourier_range_of_rows(irow, irow + 1)) {
        std::array<double, 3> kvec;  double kmag;
        grid.get_fourier_wavevector_and_norm_by_index(i, kvec, kmag);
        grid.set_fourier_from_index(i, grid.get_fourier_from_index(i) * W(kmag / Boxsize));
    }
grid.fftw_c2r();
grid.communicate_boundaries();                       // fill the extra slices
\end{cppcode}

\paragraph*{Plans and wisdom.} FFTW plans are cached and reused by every grid with the same dimension, size, alignment, thread
count and planner flag. The flag is \code{FFTW_ESTIMATE} by default; \code{set_fftw_planner_flag} (or the name
\code{estimate}, \code{measure}, \code{patient}, \code{exhaustive}) selects slower planning for faster transforms, with the grid
data preserved during planning. \code{save_fftw_wisdom} and \code{load_fftw_wisdom} store the plans so that a later run does not
measure again (valid for the same FFTW version, precision, machine and threads; collective, with task 0 doing the file I/O). The
cells of the right extra slices that FFTW may use as scratch space in an in-place transform are backed up and restored.

\paragraph*{Checkpoints and memory.} \code{dump_to_file} and \code{load_from_file} write a raw copy per task (read back with the
same decomposition; truncated files and mismatched decompositions abort with a clear message). The memory of a grid can be
computed without allocating it, and with memory logging every grid carries a label.

\paragraph*{Complex grids.} \cls{CFFTWGrid} is a complex field with complex-to-complex transforms (\code{fftw_forward},
normalised by $N_{\rm mesh}^{-D}$, and \code{fftw_backward}) and the same constructor, slabs, ghosts, index functions and
checkpoints, but no padding and all modes ($n_i\in(-N/2,N/2]$ on every axis, multiplicity 1). It holds the wave function of
the Schrödinger--Poisson solver (\cref{ch:hydro-wavedm}) and the fine grids of the NUFFT (\cref{ch:mass-assignment}).
\cls{CPencilGrid} is its pencil counterpart.

\section{The grid interface}\label{sec:grids-interface}

\file{GridInterface.h} defines the interface that generic code is written against, satisfied by \cls{FFTWGrid} and
\cls{PencilGrid} (checked at compile time by \code{has_fourier_grid_interface}). The rules:
\begin{itemize}
\item loop over Fourier modes by rows (\code{get_local_nrows_fourier}, \code{get_fourier_range_of_rows}) or with
      \code{for_each_nonzero_fourier_mode}, which passes the index, $\bm k$, $|\bm k|$ and the weight of the mode;
\item never decode a Fourier index or read the raw Fourier coordinates (local $k_x$ for slabs, global $k_x$ and local $k_y$,
      $k_z$ for pencils): use \code{get_fourier_integer_wavevector_from_index} or the wave-vector functions;
\item make grids with \code{make_grid<Grid>(Nmesh, left, right, domain)} or \code{make_grid_like} on the particle domain,
      not with a constructor; the local cells are \code{get_local_range(axis)}, and \code{CellOwners} gives the task of any cell;
\item never reuse a real-space index between grids with different ghost layers (a pencil index depends on them).
\end{itemize}
Typical generic code is templated on the grid type with a fast path for \cls{FFTWGrid} behind \code{if constexpr
(is_fftwgrid_v<...>)}.

\section{Fourier resampling}\label{sec:grids-resampling}

\code{resize_fourier_grid(input, nmesh)} copies the Fourier coefficients to a grid of another size, without interpolation or
volume factors. A mode survives only if every signed component satisfies $2|n_i|<\min(N_{\rm in},N_{\rm out})$: the ambiguous
Nyquist planes of even meshes are dropped, the highest modes of odd meshes kept, and all other output modes are zero. This is
the zero padding of the de-aliased LPT (\cref{ch:lpt}), the coarse multigrid solves and the change of resolution of the
initial conditions.

For repeated resampling a \cls{FourierResamplingPlan} gathers both slab layouts once, determines the routing per retained
x-slice as contiguous runs along the last axis (merged when both input and output indices are contiguous), and fixes the
largest payload of every cyclic peer step; applications send only complex values in bounded \code{MPI_Sendrecv} messages
(64\,MiB by default, the minimum over tasks). Invalid reuse or aliasing throws collectively before the output is touched.
\file{CopyFourierModes.h} does the same between grids of any size, decomposition and type (e.g.\ pencils), sending each mode to
the task that has it.

\begin{validation}[Validation: Fourier resampling]
Every owned output coefficient against an independent signed-wave oracle (excluded modes exactly zero, DC and non-zero Nyquist
input) in 2D, 3D and 4D, up, down and same size, odd and even meshes, a reused plan with new values, messages of one value and
not aligned with rows, task-dependent caps, and the preservation of the input and the halos.
\end{validation}

\section{Generic MPI grids}\label{sec:grids-mpigrid}

\cls{MPIGrid} is a grid of any type $T$ in any dimension, split in x-slabs or, given a \cls{Domain}, in x--y pencils, with any
number of ghost layers along the split axes. The main cells have the local index $((i_x N_{y,\rm loc}+i_y)N+i_z)$; the extra
x-slices are before and after the main grid in memory, while the y-layers of pencils are stored after them (including the
corners), so they are reached through \code{index_from_coord}, \code{get_neighbor_gridindex} or \code{get_shifted_index}.
Axes that are not split wrap periodically in the local grid. The index helpers are inlined because they sit in the inner loops of
the solvers. \cls{MPIMultiGrid} is the hierarchy of \cls{MPIGrid}s with half the cells per level, with restriction (the average of
the $2^D$ fine cells; the fine cells of a coarse cell are on the same task) used by the multigrid solvers (\cref{ch:multigrid}).
\file{ConvertMPIGridFFTWGrid.h} copies the main grid between an \cls{MPIGrid} and an \cls{FFTWGrid} (or any grid with the
interface, on its own domain).

\cls{GhostedCells} holds per-cell data of any trivially copyable type (labels, flags, indices) on the local cells of a slab or
pencil grid with ghost layers filled periodically from the neighbours, corners included (the higher axes are exchanged first, then
the lower ones carry their ghosts); tasks without cells take no part. It is what the cosmic-web, void and watershed codes use for
their labels.

\begin{validation}[Validation: grids]
\file{tests/test_grids.cpp} on 1, 2 and 4 tasks: FFT round trips with extra slices (with \code{FFTW_MEASURE} and the 2D case
where FFTW needs more memory than the local slices), move semantics, checkpoints (and aborts on truncated files, a different
decomposition or an unwritable path), ghost communication on thin slabs and non-periodic grids, index arithmetic and operators.
\file{tests/test_complex_grids.cpp}: the complex transforms against the exact coefficients of plane waves (both signs, Nyquist),
slabs equal to pencils. \file{tests/test_grid_generic.cpp}: the power-spectrum functions give the same results on
\cls{FFTWGrid} and \cls{PencilGrid}. \file{tests/test_ghosted_cells.cpp}: every ghost and corner holds the cell it wraps to.
\end{validation}
